\section{The xlog Language}\label{sec:lang}

This section describes the xlog surface: types, expressions, user-defined functions, modules, and stratified aggregation. Probabilistic facts (\texttt{p::f.}) and neural predicate declarations (\texttt{nn/k}) are language-level constructs covered in \Cref{sec:prob,sec:nesy}.

\subsection{Design Principles}

Three principles shape the surface, each a deliberate enabler of GPU execution rather than a concession. \textbf{Typed predicates:} every predicate has a statically known signature over a closed set of scalar types (\texttt{u32}, \texttt{u64}, \texttt{i32}, \texttt{i64}, \texttt{f32}, \texttt{f64}, \texttt{bool}, \texttt{symbol}), supplied by a \texttt{pred} declaration or inferred during compilation, so argument mismatches are caught before any kernel is generated and allocation stays bounded and columnar. \textbf{Classified dependencies:} negation, aggregation, modal operators, and recursion interact through SCC analysis on the predicate dependency graph. The ordinary deterministic core requires a monotone or stratified schedule, while supported nonmonotone probabilistic and epistemic components enter dedicated plans, including GPU-backed well-founded semantics for supported recursion through negation. \textbf{One language for many paradigms:} probabilistic facts, annotated disjunctions, neural predicates, and SAT constraints share the syntactic core of deterministic Datalog; parsing, normalization, and dependency analysis are shared, while backend-specific lowering preserves the semantics needed by each plan. Unsupported shapes and unbounded terms are rejected before execution rather than entering an implicit fallback.

A minimal program exhibits the core surface:

\lstinputlisting[style=xlogstyle,
                 caption={Minimal xlog program: typed predicates, facts, a recursive rule, a query.},
                 label={lst:hello}]{examples/01_hello.xlog}

\subsection{Type System}\label{sec:lang:types}

Every predicate is assigned a schema over the eight scalar types, with \texttt{f32}/\texttt{f64} following IEEE~754. A \texttt{pred} declaration supplies that schema explicitly; when it is omitted, compilation infers the schema from facts, rule heads, typed body columns, and arithmetic bindings. String literals (\texttt{"alice"}) map to \texttt{symbol} through a global symbol table, giving the runtime dense integer identifiers while preserving source-level readability; \texttt{domain} declarations introduce named aliases. Type consistency is enforced by predicate schemas and lowering-time checks: variable occurrences must agree with known column types, constants are validated against their destination columns, and arithmetic expressions preserve scalar types.

\lstinputlisting[style=xlogstyle,
                 caption={A well-typed xlog program: typed bridge predicate.},
                 label={lst:types-ok}]{examples/02_types_ok.xlog}

Declaring \texttt{bridge} so that its head argument is constrained at two incompatible types (e.g.\ \texttt{symbol} in the head but \texttt{u32} through \texttt{connected}) rejects the program at compile time.

\subsection{Expressions, Arithmetic, and User-Defined Functions}\label{sec:lang:arith}

Rule bodies may contain arithmetic, comparisons, and calls. Arithmetic bindings use the \texttt{is} operator (\texttt{D is X + Y}); the right-hand side supports \texttt{+ - * / \%}, the built-ins \texttt{abs}, \texttt{min}, \texttt{max}, \texttt{pow}, \texttt{cast}, and user-defined functions. Comparisons compile to \texttt{Filter} nodes and push below joins where profitable. Float comparison uses a total ordering---\texttt{NaN} and signed zeros yield deterministic results rather than contaminating downstream aggregations.

\lstinputlisting[style=xlogstyle,
                 caption={Arithmetic in rule bodies: pairwise proximity via squared Euclidean distance.},
                 label={lst:arith}]{examples/03_arith.xlog}

A user-defined function (UDF) is introduced with \texttt{func}, with optional parameter types and a return type after \texttt{->}. Arithmetic bodies (optionally with \texttt{if/then/else}) inline into the same \texttt{Expr} trees as built-in arithmetic and participate in the same predicate-pushdown and expression-level optimization. A predicate body, as in \texttt{func get\_parent(Child) = Parent :- parent(Child, Parent).}, contributes its relational literals immediately before the calling \texttt{is} binding in an ordinary rule or integrity constraint. Its relation may contribute zero, one, or many caller rows; the function syntax does not impose uniqueness, while ordinary set semantics deduplicates identical projected tuples. Each invocation alpha-renames its non-parameter result and body-local variables, and multiple calls expand from left to right. Nested expansion is depth-bounded; relational calls in conditional result branches and non-term arithmetic arguments are rejected rather than hoisted or coerced with changed semantics.

\lstinputlisting[style=xlogstyle,
                 caption={User-defined function: \texttt{clamp01} normalizes a raw score into the unit interval.},
                 label={lst:udf}]{examples/04_udf.xlog}

\subsection{Modules and Imports}\label{sec:lang:modules}

Programs decompose into \texttt{.xlog} modules on a search path. The \texttt{use} directive loads a module (\texttt{use graph.}) or selectively imports names (\texttt{use utils/math::\{abs, clamp\}.}); a \texttt{private} modifier hides declarations from importers. The resolver traverses the import graph, detects cycles, validates participating declarations, and performs bounded cross-module schema inference for undeclared predicate contributions before merging them. The resulting logical program then undergoes full compiler type inference, dependency analysis, and route-specific execution planning, so imported rules participate as if written inline.

\lstinputlisting[style=xlogstyle,
                 caption={Module \texttt{graph\_lib.xlog}: predicate declarations and recursive reach rules.},
                 label={lst:graph-lib}]{examples/graph_lib.xlog}

\lstinputlisting[style=xlogstyle,
                 caption={Entry \texttt{graph\_main.xlog}: imports \texttt{graph\_lib} and queries \texttt{reach}.},
                 label={lst:graph-main}]{examples/graph_main.xlog}

\subsection{Aggregations and Constraints}\label{sec:lang:agg}

Aggregation operators---\texttt{count}, \texttt{sum}, \texttt{min}, \texttt{max}, \texttt{logsumexp}---appear at rule heads with an explicit aggregation variable and lower to \texttt{GroupBy} nodes executed as radix-sort-and-reduce kernels; the stratifier rejects aggregation inside a recursive SCC. Integrity constraints are headless rules (\texttt{:- body.}) asserting that \texttt{body} is unsatisfiable once evaluation reaches fixpoint; a violation raises a diagnostic. \texttt{logsumexp} is included because it is the workhorse aggregation for probabilistic circuits, where log-space numerical stability is essential (\Cref{sec:prob}).

\lstinputlisting[style=xlogstyle,
                 caption={Stratified aggregation: per-source out-degree with an integrity constraint.},
                 label={lst:agg}]{examples/06_agg.xlog}

\subsection{Completeness over Finite Domains}\label{sec:lang:completeness}

The surface extends toward language completeness while preserving GPU-native execution: accepted constructs normalize into the typed AST and lower through their semantic backend, while forms that cannot be lowered safely are rejected at compile time rather than executed on a hidden CPU interpreter. The accepted extensions are finite typed lists (\texttt{[]}, \texttt{[H|T]}, the \texttt{list<T>} column type), statically-resolvable meta-predicates (\texttt{ground}, \texttt{functor}, \texttt{=..}, source-fact \texttt{findall}, \texttt{maplist}), explicit stratified negation-as-failure over deterministic relations, magic-set rewriting of bound recursive queries, and finite probabilistic aggregates. Deterministic extensions lower to relational joins and aggregations in RIR; probabilistic aggregates preserve their probabilistic intermediate representation. Open-ended generators, dynamic database mutation, unrestricted \texttt{call/N}, and runtime-variable predicate names remain outside the contract.
